%% ============================================================
%%  Survey: Memory Mechanisms and Autonomous Learning in LLMs
%%  Author: [Author Name]
%%  Date: March 2026
%% ============================================================
\documentclass[11pt,a4paper]{article}

%%--- Packages ---
%% 移除原有的 fontenc 和 inputenc，替换为 ctex 宏包以支持 XeLaTeX 下的中文编译
\usepackage[UTF8]{ctex} 

\usepackage{geometry}
\geometry{margin=2.5cm}
\usepackage{hyperref}
\hypersetup{colorlinks=true,linkcolor=blue,citecolor=blue,urlcolor=blue}
\usepackage{amsmath,amssymb}
\usepackage{graphicx}
\usepackage{booktabs}
\usepackage{multirow}
\usepackage{array}
\usepackage{xcolor}
\usepackage{listings}
\usepackage{enumitem}
\usepackage{caption}
\usepackage{subcaption}
\usepackage{natbib}
\usepackage{lmodern}
\usepackage{microtype}
\usepackage{tcolorbox}
\tcbuselibrary{skins,breakable}

%%--- Custom Colors ---
\definecolor{codeblue}{RGB}{0,90,180}
\definecolor{codegray}{RGB}{128,128,128}
\definecolor{codebg}{RGB}{248,248,248}
\definecolor{highlight}{RGB}{255,246,220}
\definecolor{boxframe}{RGB}{200,170,80}

%%--- Code Listing Style ---
\lstdefinestyle{pythonstyle}{
  backgroundcolor=\color{codebg},
  commentstyle=\color{codegray}\itshape,
  keywordstyle=\color{codeblue}\bfseries,
  stringstyle=\color{red!60!black},
  basicstyle=\ttfamily\footnotesize,
  breakatwhitespace=false,
  breaklines=true,
  keepspaces=true,
  numbers=left,
  numberstyle=\tiny\color{codegray},
  showspaces=false,
  showstringspaces=false,
  showtabs=false,
  tabsize=2,
  language=Python,
  frame=single,
  rulecolor=\color{gray!40}
}

%%--- Custom Box ---
\newtcolorbox{insightbox}[1][]{
  enhanced,
  breakable,
  colback=highlight,
  colframe=boxframe,
  title=#1,
  fonttitle=\bfseries,
  boxrule=1pt,
  arc=3pt
}

%%--- Title ---
\title{
  \vspace{-1cm}
  \LARGE\bfseries 大语言模型记忆机制与自主学习前沿调研报告\\[0.5em]
  \large Survey on Memory Mechanisms and Autonomous Learning\\
  in Large Language Models (2023--2026)
}
\author{调研执行：AI 算法研究助理\\[0.3em]
  \small 工作区：\texttt{2026Winter} $\cdot$ 日期：2026年3月}
\date{}

\begin{document}

\maketitle
\thispagestyle{empty}

\begin{abstract}
本报告系统调研了大语言模型（LLM）记忆机制与自主学习领域2023至2026年间的核心文献，
涵盖工作区已有论文（共8篇）及本次新检索下载的5篇高价值前沿论文（已存入
\texttt{2026Winter/More}目录）。报告从记忆的形成、检索、更新与整合四个维度对各方法进行
横向对比，揭示当前主流范式的核心创新与固有局限，并针对最优两篇工作——\textbf{Memory-R1}
与\textbf{Self-Consolidation}——给出具体的代码复现路线与实验框架建议，以期为后续研究提供
参考。
\end{abstract}

\tableofcontents
\newpage

%% ===========================================================
\section{研究背景与动机}
%% ===========================================================

大语言模型在推理、生成和工具调用方面展现出强大能力，但其本质上是\textbf{无状态}
（stateless）的系统：每次推理均从空白上下文出发，无法跨对话保留经验。随着应用场景扩展
至长期对话、个性化助手、多任务代理等领域，\textbf{持久记忆}（persistent memory）与
\textbf{自主学习}（autonomous learning）能力正成为制约 LLM 实际部署的核心瓶颈。

\subsection{记忆的多维度定义}

受认知科学启发，LLM 记忆研究通常借鉴人类记忆的分类体系：

\begin{itemize}
  \item \textbf{情景记忆}（Episodic Memory）：对话历史、事件序列；
  \item \textbf{语义记忆}（Semantic Memory）：提炼的知识与偏好；
  \item \textbf{程序记忆}（Procedural Memory）：内化于模型参数中的技能。
\end{itemize}

\subsection{核心挑战}

\begin{enumerate}
  \item \textbf{上下文窗口限制}：即便最大型模型也无法在推理时容纳无限历史；
  \item \textbf{记忆遗忘与更新}：随时间推移，旧记忆应被精化或淘汰，而非单纯堆积；
  \item \textbf{检索效率}：从大规模外部记忆库中快速定位相关片段；
  \item \textbf{噪声与冗余}：朴素的累积策略会引入噪声，降低记忆质量；
  \item \textbf{参数化整合}：如何将外部经验真正"内化"至模型权重。
\end{enumerate}

%% ===========================================================
\section{调研论文全景}
%% ===========================================================

\subsection{已有文献（工作区原有）}

\begin{table}[h!]
\centering
\caption{工作区已有论文概览}
\label{tab:existing}
\renewcommand{\arraystretch}{1.3}
\begin{tabular}{llllp{5.5cm}}
\toprule
\textbf{论文简称} & \textbf{年份} & \textbf{存放位置} & \textbf{链接} & \textbf{核心贡献} \\
\midrule
Generative Agents & 2023 & TTL & \href{https://arxiv.org/abs/2304.03442}{ArXiv} & 模拟人类行为的 Agent，
  引入记忆流（memory stream）、记忆检索与反思机制 \\
Reflexion & 2023 & TTL & \href{https://arxiv.org/abs/2303.11366}{ArXiv} & 语言强化学习，通过语言形式的自我反思作为记忆信号 \\
MemGPT & 2023 & Self-evolving & \href{https://arxiv.org/abs/2310.08560}{ArXiv} & 仿 OS 分层记忆管理，主内存+外部存储+分页机制 \\
MemoryBank & 2023 & Self-evolving & \href{https://arxiv.org/abs/2305.10250}{ArXiv} & 记忆银行框架，Ebbinghaus 遗忘曲线启发的记忆更新 \\
RAG & 2020 & Self-evolving & \href{https://arxiv.org/abs/2005.11401}{ArXiv} & 检索增强生成，将外部知识库实时注入上下文 \\
Self-RAG & 2023 & Self-evolving & \href{https://arxiv.org/abs/2310.11511}{ArXiv} & 自我批评的检索增强，引入 Critique Token 控制检索时机 \\
Evo-Memory & 2025 & 根目录 & \href{https://arxiv.org/abs/2511.20857}{ArXiv} & 测试时自演化记忆基准，提出 ReMem 流水线 \\
Conditional Memory & 2026 & 根目录 & \href{https://github.com/deepseek-ai/Engram}{GitHub} & 条件稀疏记忆，KV Cache 新维度的结构化稀疏 \\
\bottomrule
\end{tabular}
\end{table}

\subsection{新下载文献（存入 \texttt{More/} 目录）}

\begin{table}[h!]
\centering
\caption{新检索高价值论文概览}
\label{tab:new}
\renewcommand{\arraystretch}{1.3}
\begin{tabular}{lllll}
\toprule
\textbf{文件名} & \textbf{ArXiv ID} & \textbf{年份} & \textbf{链接} & \textbf{核心方向} \\
\midrule
2025-Xu-A-Mem.pdf & 2502.12110 & 2025 & \href{https://arxiv.org/abs/2502.12110}{ArXiv} & Zettelkasten 启发的主动记忆组织 \\
2025-Chhikara-Mem0.pdf & 2504.19413 & 2025 & \href{https://arxiv.org/abs/2504.19413}{ArXiv} & 生产级可扩展长期记忆 \\
2025-Chen-Memory-R1.pdf & 2508.19828 & 2025 & \href{https://arxiv.org/abs/2508.19828}{ArXiv} & RL 驱动的记忆主动管理 \\
2025-Wang-PAMU.pdf & 2510.09720 & 2025 & \href{https://arxiv.org/abs/2510.09720}{ArXiv} & 偏好感知的记忆动态更新 \\
2026-Zhang-Self-Consolidation.pdf & 2602.01966 & 2026 & \href{https://arxiv.org/abs/2602.01966}{ArXiv} & 文本经验到参数的自我整合 \\
\bottomrule
\end{tabular}
\end{table}

%% ===========================================================
\section{各论文核心创新点详解}
%% ===========================================================

\subsection{Generative Agents}

Park 等人构建了一个包含25个 Agent 的沙盒世界模拟，首次在 LLM Agent 中引入完整的
\textbf{记忆流}（Memory Stream）。其记忆机制由三部分构成：
\begin{itemize}
  \item \textbf{记忆存储}：时间戳+事件描述的扁平序列；
  \item \textbf{记忆检索}：综合近因（recency）、重要性（importance）和相关性（relevance）的加权评分；
  \item \textbf{记忆反思}（Reflection）：定期将低层记忆归纳为高层洞见。
\end{itemize}

该工作开创性地将\emph{反思}（reflection）作为记忆压缩手段，但缺乏参数化更新，记忆仅
以文本形式外存。

\subsection{Reflexion}

Reflexion 将错误经验转化为\textbf{语言形式的自我反思文本}，存入专属的 Episodic Memory
Buffer，在后续任务中作为 few-shot 示例提供给模型。其本质是一种\textbf{语言强化学习}
（verbal reinforcement learning），极大降低了梯度更新的依赖。局限在于记忆只保存成功轨迹
的反思摘要，缺乏对失败模式的系统性总结与利用。

\subsection{MemGPT}

MemGPT 将操作系统的\textbf{分页内存管理}理念迁移到 LLM：上下文窗口视为"主内存"，
外部数据库视为"磁盘"；模型通过特殊函数调用（\texttt{memory.append}、
\texttt{memory.search}）主动控制信息的"换入换出"。这是首个将 LLM 自身作为内存
控制器（memory controller）的工作，但内存操作仍依赖人工设计的函数接口，缺乏学习性。

\subsection{MemoryBank}

MemoryBank 受 Ebbinghaus 遗忘曲线启发，为每条记忆引入\textbf{衰减函数}：
$$
\text{retention}(t) = e^{-\frac{t}{\tau}}
$$
近期交互的记忆保留率高，遥远记忆自然衰减。更新时通过 LLM 对旧记忆摘要进行改写整合。
虽然模拟了遗忘机制，但参数 $\tau$ 固定，无法自适应不同用户行为节律。

\subsection{RAG \& Self-RAG}

RAG 通过稠密检索将外部知识实时注入上下文，解决了参数知识滞后问题，但对话历史管理
的颗粒度较粗。Self-RAG 引入 Critique Token（\texttt{[Retrieve]}、\texttt{[IsRel]}、
\texttt{[IsSup]}）让模型自主决策是否检索及检索结果质量，显著提升效率；然而其记忆系统
本质仍是静态语料库，不支持在线更新。

\subsection{Evo-Memory}

Evo-Memory 提出\textbf{测试时演化记忆}（test-time evolving memory）基准，评估 Agent
在顺序任务流中持续检索、整合和更新记忆的能力。配套提出的 \textbf{ReMem} 框架（Action
$\to$ Think $\to$ Memory Refine 流水线）以迭代记忆精化取代朴素累积，是首个针对记忆演化
质量设计专项基准的工作。其不足在于记忆仍以文本形式外存，参数化能力有限。

\subsection{Conditional Memory}

Conditional Memory 从\textbf{模型结构}视角切入，将 KV Cache 的稀疏化从"位置维度"
扩展到"内容条件维度"，在大规模 LLM 推理中动态路由相关 KV 对。这本质上是一种
\emph{参数内记忆}（in-weights / in-cache memory）的高效压缩方案，与外部记忆研究形成
互补，但不直接解决跨会话的长期记忆问题。

\subsection{A-MEM}

A-MEM 将卡片笔记法（\textbf{Zettelkasten}）引入 LLM 记忆：每条新记忆被构建为
一张"知识卡片"，包含上下文描述、关键词、标签及指向历史记忆的显式链接；当新记忆
加入时，相关旧记忆的属性也被\emph{主动更新}（memory evolution）。与 RAG 的被动检索
相比，A-MEM 实现了\textbf{主动记忆组织}（agentic memory organization），在六个基础
模型上均超越 SOTA 基线，并公开了代码。

\subsection{Mem0}

Mem0 面向\textbf{生产部署}需求，构建了可扩展的记忆中心架构：LLM 动态提取并
整合对话中的关键信息，图结构记忆模块进一步捕捉对话元素间的关系网络。在 LOCOMO
基准上，Mem0 相比全上下文方法将 p95 延迟降低 \textbf{91\%}、token 成本节省超
\textbf{90\%}，同时在 LLM-as-a-Judge 指标上超越 OpenAI 26\%。其核心贡献是
\textbf{工程效率}与\textbf{准确性的同步提升}。

\subsection{Memory-R1}

Memory-R1 用\textbf{强化学习}（RL）替代启发式规则来驱动记忆管理，包含两个专用
Agent：
\begin{itemize}
  \item \textbf{Memory Manager}：通过 PPO/GRPO 学习 ADD、UPDATE、DELETE、NOOP 四种
    操作策略；
  \item \textbf{Answer Agent}：预选相关记忆条目并推理作答。
\end{itemize}
关键特性：\textbf{仅需 152 条训练 QA 对}即可在 LoCoMo、MSC、LongMemEval 三个基准
上超越强基线，且泛化至 3B--14B 多种模型规模，是目前记忆管理中\textbf{数据效率最高}
的方案之一。

\subsection{PAMU}

PAMU 专注于\textbf{记忆更新}这一被忽视的环节，提出用滑动窗口均值（SW）与指数移动
均值（EMA）的融合来构建偏好感知表示：
$$
\text{PAMU}(t) = \alpha \cdot \overline{m}_{\text{SW}}(t) + (1-\alpha) \cdot \overline{m}_{\text{EMA}}(t)
$$
SW 捕捉短期行为波动，EMA 保留长期用户偏好，两者互补。在 LoCoMo 数据集的五个任务
场景中，PAMU 可无缝插入5种基线系统并带来一致提升，体现了良好的\textbf{模块化设计}。

\subsection{Self-Consolidation}

Self-Consolidation 是迄今最前沿的自演化 Agent 工作（2026年2月），提出双轨演化机制：
\begin{enumerate}
  \item \textbf{对比反思}（Contrastive Reflection）：对比成功与失败轨迹，显式归纳
    易错模式，而非仅存储成功案例；
  \item \textbf{自我整合}（Self-Consolidation）：将文本形式的历史经验通过低秩微调
    （LoRA）蒸馏进模型参数，避免外部记忆库随时间膨胀的问题。
\end{enumerate}
这是首个将\textbf{参数化学习}与\textbf{外部记忆检索}统一到单一自演化框架的工作，
从根本上突破了"记忆只能在上下文中，不能改变模型本身"的范式边界。

%% ===========================================================
\section{记忆更新机制横向对比}
%% ===========================================================

\begin{table}[h!]
\centering
\caption{各方法记忆更新机制对比}
\label{tab:compare}
\renewcommand{\arraystretch}{1.4}
\small
\begin{tabular}{p{3.2cm}cccccc}
\toprule
\textbf{方法} &
\textbf{外部存储} &
\textbf{参数更新} &
\textbf{主动更新} &
\textbf{遗忘/衰减} &
\textbf{失败反思} &
\textbf{在线学习} \\
\midrule
Generative Agents & \checkmark & $\times$ & 半自动 & $\times$ & $\times$ & $\times$ \\
Reflexion          & \checkmark & $\times$ & 手动触发 & $\times$ & 部分 & $\times$ \\
MemGPT             & \checkmark & $\times$ & 函数调用 & $\times$ & $\times$ & $\times$ \\
MemoryBank         & \checkmark & $\times$ & 改写整合 & Ebbinghaus & $\times$ & $\times$ \\
RAG                & \checkmark & $\times$ & $\times$ & $\times$ & $\times$ & $\times$ \\
Self-RAG           & \checkmark & $\times$ & 自主决策 & $\times$ & $\times$ & $\times$ \\
Evo-Memory         & \checkmark & $\times$ & 迭代精化 & $\times$ & $\times$ & 测试时 \\
A-MEM               & \checkmark & $\times$ & 主动演化 & $\times$ & $\times$ & 测试时 \\
Mem0               & \checkmark & $\times$ & 自动提取 & 隐式 & $\times$ & \checkmark \\
PAMU               & \checkmark & $\times$ & SW+EMA & 加权衰减 & $\times$ & \checkmark \\
Memory-R1          & \checkmark & $\times$ & RL 学习 & RL 策略 & $\times$ & \checkmark \\
Self-Consolidation & \checkmark & \textbf{\checkmark} & RL+对比 & LoRA 蒸馏 & \checkmark & \checkmark \\
Conditional Memory & $\times$ & \checkmark & KV 稀疏 & 动态路由 & $\times$ & $\times$ \\
\bottomrule
\end{tabular}
\end{table}

\subsection{维度深析}

\paragraph{记忆形成（Formation）}
从 Generative Agents 的"事件日志堆积"到 A-MEM 的"结构化知识卡片"，再到
Self-Consolidation 的"参数蒸馏"，记忆形成经历了从\emph{被动记录}→\emph{主动组织}
→\emph{内化参数}的三代演进。

\paragraph{记忆检索（Retrieval）}
RAG 族方法基于稠密相似度；MemGPT 依赖显式函数签名；Memory-R1 的 Answer Agent
通过 RL 学得"何时检索、检索什么"的策略，是首个将检索本身视为\emph{可学习决策}的工作。

\paragraph{记忆更新（Update）}
PAMU 揭示现有方法在更新阶段的普遍疏忽，SW+EMA 融合是一种\textbf{低成本插件式}
解决方案。Memory-R1 的 RL 框架则端到端学习更新策略，但训练成本更高。

\paragraph{记忆整合（Consolidation）}
只有 Self-Consolidation 真正实现了文本记忆→模型参数的整合，突破了外部记忆随时间
无限膨胀的瓶颈。这与神经科学中"系统性记忆巩固"（system consolidation）高度吻合。

%% ===========================================================
\section{现有方法的局限性}
%% ===========================================================

\begin{enumerate}
  \item \textbf{记忆更新被低估}：绝大多数工作将研究重心放在检索上，而忽视旧记忆的
    动态精化。当用户偏好随时间发生漂移时，静态记忆会导致推荐/回复质量明显下滑。
    PAMU 仅部分解决此问题，仍缺乏对\emph{语义级}记忆冲突的检测和消解机制。

  \item \textbf{失败经验利用不足}：Reflexion 等方法仅保存成功轨迹的反思，忽略失败
    模式的教学价值。Self-Consolidation 是首个系统性挖掘失败案例的工作，表明负样本
    对长期演化至关重要。

  \item \textbf{参数化整合缺失}：外部记忆库本质上是"追加型数据库"，随时间增长必然
    带来检索延迟上升和噪声累积。只有 Self-Consolidation 迈出了参数整合的第一步，
    但其 LoRA 微调频率与触发时机的设计仍较启发式，缺乏理论依据。

  \item \textbf{评估基准不统一}：不同方法在 LoCoMo、MSC、LongMemEval 等数据集上
    各自报告，交叉对比困难。Evo-Memory 基准是一次有益尝试，但覆盖范围有限。

  \item \textbf{多 Agent 场景稀缺}：现有方法几乎全部聚焦单 Agent，共享记忆池、
    记忆所有权、一致性等多 Agent 问题几乎无人涉及。

  \item \textbf{隐私与安全}：持久外部记忆本质上存储了用户敏感交互，记忆的权限管理、
    差分隐私保护、记忆"遗忘权"等议题亟待研究。
\end{enumerate}

%% ===========================================================
\section{顶级论文代码复现思路}
%% ===========================================================

\subsection{Memory-R1：基于强化学习的记忆管理}

\subsubsection{核心架构设计}

Memory-R1 包含两个 RL 微调的 Agent，建议采用如下实验框架：

\begin{lstlisting}[style=pythonstyle, caption={Memory-R1 核心组件伪代码}]
# ============================================================
# Memory-R1 Reproduction Framework
# 推荐环境: Python 3.10+, PyTorch 2.x, transformers, trl
# ============================================================

from dataclasses import dataclass, field
from typing import Literal
import torch
from transformers import AutoModelForCausalLM, AutoTokenizer
from trl import PPOTrainer, PPOConfig

MemoryOp = Literal["ADD", "UPDATE", "DELETE", "NOOP"]

@dataclass
class MemoryEntry:
    """外部记忆库条目"""
    content: str
    timestamp: float
    relevance_score: float = 0.0

class MemoryStore:
    """外部记忆数据库（可替换为 FAISS / ChromaDB）"""
    def __init__(self, max_entries: int = 500):
        self.entries: list[MemoryEntry] = []
        self.max_entries = max_entries

    def add(self, content: str, timestamp: float):
        self.entries.append(MemoryEntry(content=content, timestamp=timestamp))
        if len(self.entries) > self.max_entries:
            self.entries.pop(0)  # FIFO truncation

    def retrieve(self, query: str, topk: int = 5) -> list[MemoryEntry]:
        # 建议替换为 Dense Retrieval (DPR / BM25)
        return sorted(self.entries, key=lambda e: e.relevance_score,
                      reverse=True)[:topk]

class MemoryManager:
    """Memory Manager Agent - PPO/GRPO 微调目标"""
    def __init__(self, model_name: str, store: MemoryStore):
        self.model = AutoModelForCausalLM.from_pretrained(model_name)
        self.tokenizer = AutoTokenizer.from_pretrained(model_name)
        self.store = store

    def predict_operation(self, new_event: str,
                          history: list[MemoryEntry]
                          ) -> tuple[MemoryOp, str]:
        """输出操作类型 + 操作目标内容"""
        prompt = self._build_manager_prompt(new_event, history)
        output = self._generate(prompt)
        op, content = self._parse_output(output)
        return op, content

    def _build_manager_prompt(self, event: str,
                              history: list[MemoryEntry]) -> str:
        memory_str = "\n".join(
            [f"[{i}] {e.content}" for i, e in enumerate(history)])
        return (
            f"Current memory entries:\n{memory_str}\n\n"
            f"New event: {event}\n\n"
            f"Decide: ADD / UPDATE [index] / DELETE [index] / NOOP\n"
            f"Operation:"
        )

class AnswerAgent:
    """Answer Agent - 选择相关记忆并作答"""
    def __init__(self, model_name: str, store: MemoryStore):
        self.model = AutoModelForCausalLM.from_pretrained(model_name)
        self.tokenizer = AutoTokenizer.from_pretrained(model_name)
        self.store = store

    def answer(self, question: str) -> str:
        relevant = self.store.retrieve(question, topk=5)
        prompt = self._build_answer_prompt(question, relevant)
        return self._generate(prompt)

# RL 训练配置
def train_memory_r1(manager: MemoryManager, data: list[dict]):
    """
    data 格式: [{"events": [...], "question": str, "answer": str}, ...]
    奖励函数: 最终答案 F1/Exact Match
    """
    config = PPOConfig(
        model_name=manager.model.name_or_path,
        learning_rate=1.41e-5,
        batch_size=16,
        mini_batch_size=4,
        gradient_accumulation_steps=2,
        ppo_epochs=4,
    )
    # outcome-driven reward: 仅基于最终 QA 准确率
    def reward_fn(completions, answers):
        return [1.0 if c.strip() == a.strip() else 0.0
                for c, a in zip(completions, answers)]
    # ... PPO training loop
\end{lstlisting}

\subsubsection{实验框架与建议}

\begin{insightbox}[Memory-R1 复现实验框架]
\textbf{数据集}：LoCoMo（首选）、MSC、LongMemEval \\
\textbf{基础模型}：Qwen2.5-3B-Instruct（快速验证）→ Qwen2.5-7B/14B（正式实验） \\
\textbf{训练规模}：仅需 $\sim$152 QA 对，建议设置消融实验：50/100/152/300 \\
\textbf{检索模块}：初步用 BM25，进阶替换 DPR + FAISS 向量库 \\
\textbf{评估指标}：F1, Exact Match, LLM-as-a-Judge（GPT-4o 评分） \\
\textbf{对比基线}：MemGPT, MemoryBank, Mem0, A-MEM \\
\textbf{关键消融}：(1) 仅 ADD 操作；(2) 无 RL，用贪婪策略；(3) 无 Answer Agent
\end{insightbox}

\begin{lstlisting}[style=pythonstyle, caption={完整实验运行脚本骨架}]
# run_memoryr1_exp.py
import argparse
from datasets import load_dataset

def main(args):
    # 1. 加载数据
    dataset = load_dataset("locomo_bench")  # 替换为实际路径

    # 2. 初始化组件
    store = MemoryStore(max_entries=args.max_memory)
    manager = MemoryManager(args.model_name, store)
    agent = AnswerAgent(args.model_name, store)

    # 3. Phase 1: SFT 热启动 (可选，提升训练稳定性)
    # sft_train(manager, agent, dataset["train"], epochs=1)

    # 4. Phase 2: RL 微调
    train_memory_r1(manager, dataset["train"])

    # 5. 评估
    results = evaluate(agent, dataset["test"])
    print(f"F1: {results['f1']:.4f} | EM: {results['em']:.4f}")

if __name__ == "__main__":
    parser = argparse.ArgumentParser()
    parser.add_argument("--model_name", default="Qwen/Qwen2.5-7B-Instruct")
    parser.add_argument("--max_memory", type=int, default=200)
    args = parser.parse_args()
    main(args)
\end{lstlisting}

\subsection{Self-Consolidation：文本经验的参数内化}

\subsubsection{核心架构设计}

Self-Consolidation 的关键在于"对比反思"与"LoRA 蒸馏"的协同，复现重点在于正确
构建对比对（contrastive pair）和设计合理的 LoRA 触发策略。

\begin{lstlisting}[style=pythonstyle, caption={Self-Consolidation 核心框架}]
# ============================================================
# Self-Consolidation Reproduction Framework
# 依赖: peft (LoRA), transformers, torch
# ============================================================

from peft import LoraConfig, get_peft_model, TaskType
import torch.nn.functional as F

@dataclass
class Trajectory:
    """单次任务执行轨迹"""
    steps: list[str]          # 行动序列
    success: bool             # 是否成功
    final_response: str       # 最终输出

class ContrastiveReflector:
    """
    对比反思模块：对比成功/失败轨迹，提炼错误模式与可复用洞见
    """
    def __init__(self, llm_model, tokenizer):
        self.model = llm_model
        self.tokenizer = tokenizer

    def reflect(self,
                success_traj: Trajectory,
                failure_traj: Trajectory) -> str:
        """
        输入：一对成功/失败轨迹
        输出：自然语言反思摘要（含错误模式 + 改进策略）
        """
        prompt = self._build_contrast_prompt(success_traj, failure_traj)
        reflection = self._generate(prompt, max_new_tokens=256)
        return reflection

    def _build_contrast_prompt(self, s: Trajectory, f: Trajectory) -> str:
        return (
            "Compare the following two task trajectories:\n\n"
            f"[SUCCESS]\n{chr(10).join(s.steps)}\n\n"
            f"[FAILURE]\n{chr(10).join(f.steps)}\n\n"
            "Identify:\n"
            "1. Key error patterns in the failure trajectory\n"
            "2. Reusable strategies from the successful trajectory\n"
            "3. Concrete advice to avoid future failures\n\n"
            "Reflection:"
        )

class SelfConsolidator:
    """
    自我整合模块：将积累的文本反思蒸馏为 LoRA 参数增量
    触发条件建议：每积累 N=50 条反思，或任务成功率连续下降3轮
    """
    def __init__(self, base_model, tokenizer, lora_rank: int = 16):
        lora_config = LoraConfig(
            task_type=TaskType.CAUSAL_LM,
            r=lora_rank,
            lora_alpha=32,
            target_modules=["q_proj", "v_proj"],  # 可调整目标层
            lora_dropout=0.05,
            bias="none",
        )
        self.model = get_peft_model(base_model, lora_config)
        self.tokenizer = tokenizer
        self.reflection_buffer: list[str] = []  # 文本记忆缓冲区

    def accumulate(self, reflection: str):
        """将反思文本加入缓冲区"""
        self.reflection_buffer.append(reflection)

    def consolidate(self, trigger_threshold: int = 50):
        """
        当缓冲区满足阈值时执行参数蒸馏
        采用 SFT Loss 在反思文本上微调 LoRA 层
        """
        if len(self.reflection_buffer) < trigger_threshold:
            return  # 未达到触发条件

        # 构建蒸馏数据集
        texts = self.reflection_buffer.copy()
        self.reflection_buffer.clear()

        self.model.train()
        optimizer = torch.optim.AdamW(
            self.model.parameters(), lr=1e-4, weight_decay=0.01)

        for text in texts:
            inputs = self.tokenizer(
                text, return_tensors="pt", truncation=True,
                max_length=512).to(self.model.device)
            outputs = self.model(**inputs, labels=inputs["input_ids"])
            loss = outputs.loss
            loss.backward()
            optimizer.step()
            optimizer.zero_grad()

        print(f"[SelfConsolidation] Consolidated {len(texts)} reflections"
              f" into LoRA parameters.")

class SelfEvolvingAgent:
    """
    集成 ContrastiveReflector + SelfConsolidator 的完整自演化 Agent
    """
    def __init__(self, model, tokenizer):
        self.model = model
        self.tokenizer = tokenizer
        self.reflector = ContrastiveReflector(model, tokenizer)
        self.consolidator = SelfConsolidator(model, tokenizer)
        self.trajectory_history: list[Trajectory] = []

    def execute_task(self, task: str) -> Trajectory:
        """执行任务并记录轨迹（由子类实现具体逻辑）"""
        raise NotImplementedError

    def evolve(self):
        """每完成一批任务后触发演化流程"""
        successes = [t for t in self.trajectory_history if t.success]
        failures = [t for t in self.trajectory_history if not t.success]

        # 对比反思（要求至少各有一条）
        for s, f in zip(successes, failures):
            reflection = self.reflector.reflect(s, f)
            self.consolidator.accumulate(reflection)

        # 尝试参数整合
        self.consolidator.consolidate()
        self.trajectory_history.clear()
\end{lstlisting}

\subsubsection{实验框架与建议}

\begin{insightbox}[Self-Consolidation 复现实验框架]
\textbf{数据集}：WebArena / ALFWorld（Agent 任务）+ LoCoMo（对话记忆任务） \\
\textbf{基础模型}：LLaMA-3.1-8B-Instruct 或 Qwen2.5-7B-Instruct \\
\textbf{LoRA 配置}：rank=16, alpha=32, target=\texttt{q\_proj, v\_proj} \\
\textbf{触发策略}：实验三种阈值——每 20/50/100 条反思触发一次整合 \\
\textbf{评估重点}：长期演化曲线（任务成功率 vs. 轮次），非单次测试表现 \\
\textbf{关键消融}：\\
  \quad (1) 仅对比反思，不做参数整合（验证整合的必要性）\\
  \quad (2) 仅保存成功轨迹（复现 Reflexion 基线，验证对比的价值）\\
  \quad (3) 全模型微调 vs. LoRA（效率与效果权衡）\\
\textbf{可视化}：记忆库大小随时间变化曲线；LoRA 参数范数的演化图
\end{insightbox}

%% ===========================================================
\section{研究展望}
%% ===========================================================

基于上述分析，我们建议以下几个有价值的未来研究方向：

\begin{enumerate}
  \item \textbf{记忆更新的理论框架}：将 PAMU 的 SW+EMA 融合推广至基于贝叶斯
    更新的记忆精化，为"何时遗忘、何时保留"提供理论依据。

  \item \textbf{联合训练}：Memory-R1 的 Memory Manager 与 Self-Consolidation 的
    对比反思可以端到端联合优化，探索 RL + LoRA 协同训练的可行性。

  \item \textbf{多 Agent 共享记忆}：将 Mem0 的图结构记忆扩展到多 Agent 协作场景，
    研究记忆所有权、一致性与冲突消解。

  \item \textbf{记忆可解释性}：对比 A-MEM 的显式链接网络与 Self-Consolidation 的
    隐式参数整合，探索可解释性与能力之间的权衡。

  \item \textbf{安全记忆}：引入差分隐私（DP-SGD）对记忆蒸馏过程进行保护，研究
    "记忆遗忘"（memory unlearning）在医疗、法律等高隐私场景的应用。
\end{enumerate}

%% ===========================================================
\section{总结}
%% ===========================================================

本报告系统梳理了 LLM 记忆机制从 2020 年到 2026 年初的演化脉络：从 RAG 的
\emph{被动检索}，到 MemGPT 的\emph{主动管理}，到 Memory-R1 的\emph{强化学习策略}，
再到 Self-Consolidation 的\emph{参数内化}，每一代方法都在突破上一代的核心瓶颈。

\begin{center}
\begin{tabular}{lcl}
\toprule
\textbf{代际} & \textbf{代表作} & \textbf{核心突破} \\
\midrule
第一代 & RAG, MemGPT & 外部记忆存储与基础检索 \\
第二代 & Reflexion, MemoryBank & 记忆反思与衰减机制 \\
第三代 & A-MEM, Mem0, PAMU & 主动组织与动态更新 \\
第四代 & Memory-R1, Self-Consolidation & RL 策略学习 + 参数整合 \\
\bottomrule
\end{tabular}
\end{center}

\vspace{1em}
当前最值得关注的工作是 \textbf{Memory-R1}（以极低数据成本实现高质量记忆管理）和
\textbf{Self-Consolidation}（首次将文本记忆内化为模型参数，开创第四代范式）。
建议优先复现这两篇工作，并探索二者的融合路径。

%% ===========================================================
%% References
%% ===========================================================
\newpage
\section*{参考文献}
\bibliographystyle{unsrt}

\begin{thebibliography}{99}

\bibitem{park2023generative}
Park, J.~S., et al.
\newblock Generative agents: Interactive simulacra of human behavior.
\newblock In \textit{Proceedings of the 36th Annual ACM Symposium on User
  Interface Software and Technology (UIST)}, 2023.

\bibitem{shinn2023reflexion}
Shinn, N., et al.
\newblock Reflexion: Language agents with verbal reinforcement learning.
\newblock In \textit{Advances in Neural Information Processing Systems
  (NeurIPS)}, 2023.

\bibitem{packer2023memgpt}
Packer, C., et al.
\newblock MemGPT: Towards LLMs as operating systems.
\newblock \textit{arXiv preprint}, arXiv:2310.08560, 2023.

\bibitem{zhong2024memorybank}
Zhong, W., et al.
\newblock MemoryBank: Enhancing large language models with long-term memory.
\newblock In \textit{Proceedings of AAAI}, 2024.

\bibitem{lewis2020retrieval}
Lewis, P., et al.
\newblock Retrieval-augmented generation for knowledge-intensive NLP tasks.
\newblock In \textit{Advances in Neural Information Processing Systems
  (NeurIPS)}, 2020.

\bibitem{asai2023selfrag}
Asai, A., et al.
\newblock Self-RAG: Learning to retrieve, generate, and critique through
  self-reflection.
\newblock In \textit{International Conference on Learning Representations
  (ICLR)}, 2024.

\bibitem{evomemory2025}
Anonymous.
\newblock Evo-Memory: Benchmarking LLM agent test-time learning with
  self-evolving memory.
\newblock \textit{arXiv preprint}, 2025.

\bibitem{condmem2026}
Anonymous.
\newblock Conditional memory via scalable lookup: A new axis of sparsity for
  large language models.
\newblock \textit{arXiv preprint}, 2026.

\bibitem{xu2025amem}
Xu, W., et al.
\newblock A-MEM: Agentic memory for LLM agents.
\newblock \textit{arXiv preprint}, arXiv:2502.12110, 2025.

\bibitem{chhikara2025mem0}
Chhikara, P., et al.
\newblock Mem0: Building production-ready AI agents with scalable long-term
  memory.
\newblock \textit{arXiv preprint}, arXiv:2504.19413, 2025.

\bibitem{chen2025memoryr1}
Chen, J., et al.
\newblock Memory-R1: Enhancing large language model agents to manage and
  utilize memories via reinforcement learning.
\newblock \textit{arXiv preprint}, arXiv:2508.19828, 2025.

\bibitem{wang2025pamu}
Wang, Z., et al.
\newblock Preference-aware memory update for long-term LLM agents.
\newblock \textit{arXiv preprint}, arXiv:2510.09720, 2025.

\bibitem{zhang2026selfcon}
Zhang, H., et al.
\newblock Self-consolidation for self-evolving agents.
\newblock \textit{arXiv preprint}, arXiv:2602.01966, 2026.

\end{thebibliography}

\end{document}